
\documentclass{article}
\usepackage{colortbl}
\usepackage{makecell}
\usepackage{multirow}
\usepackage{supertabular}

\begin{document}

\newcounter{utterance}

\twocolumn

{ \footnotesize  \setcounter{utterance}{1}
\setlength{\tabcolsep}{0pt}
\begin{supertabular}{c@{$\;$}|p{.15\linewidth}@{}p{.15\linewidth}p{.15\linewidth}p{.15\linewidth}p{.15\linewidth}p{.15\linewidth}}

    \# & $\;$A & \multicolumn{4}{c}{Game Master} & $\;\:$B\\
    \hline 

    \theutterance \stepcounter{utterance}  

    & & \multicolumn{4}{p{0.6\linewidth}}{\cellcolor[rgb]{0.9,0.9,0.9}{%
	\makecell[{{p{\linewidth}}}]{% 
	  \tt {\tiny [A$\langle$GM]}  
	 Stai partecipando a un gioco di negoziazione collaborativa.\\ \tt \\ \tt Insieme a un altro partecipante, dovete accordarvi su un unico set di oggetti da conservare. Ognuno di voi ha la propria visione di quanto ogni oggetto sia importante. Non sai come l'altro partecipante valuta gli oggetti. Inoltre, ti viene fornito lo sforzo che ogni oggetto richiede.\\ \tt Potete accordarvi su un set solo se lo sforzo totale degli oggetti selezionati non supera un limite condiviso:\\ \tt \\ \tt LIMITE = 9322\\ \tt \\ \tt Ecco i valori di sforzo dei singoli oggetti:\\ \tt \\ \tt Sforzo oggetto = {"A08": 380, "B95": 448, "A39": 568, "A65": 351, "B81": 896, "B34": 821, "A28": 526, "C82": 502, "B69": 326, "A87": 962, "B64": 619, "B15": 115, "C79": 598, "B16": 662, "C53": 300, "C83": 830, "C10": 562, "A97": 679, "A25": 283, "C84": 440, "A71": 12, "B23": 863, "B96": 319, "C51": 771, "C17": 89, "B19": 655, "C72": 504, "A94": 118, "A46": 514, "B43": 226, "B37": 892, "A88": 621, "A96": 766, "C60": 659, "A75": 768}\\ \tt \\ \tt Ecco la tua visione personale sull'importanza di ogni oggetto:\\ \tt \\ \tt Valori di importanza dell'oggetto = {"A08": 476, "B95": 544, "A39": 664, "A65": 447, "B81": 992, "B34": 917, "A28": 622, "C82": 598, "B69": 422, "A87": 1058, "B64": 715, "B15": 211, "C79": 694, "B16": 758, "C53": 396, "C83": 926, "C10": 658, "A97": 775, "A25": 379, "C84": 536, "A71": 108, "B23": 959, "B96": 415, "C51": 867, "C17": 185, "B19": 751, "C72": 600, "A94": 214, "A46": 610, "B43": 322, "B37": 988, "A88": 717, "A96": 862, "C60": 755, "A75": 864}\\ \tt \\ \tt Obiettivo:\\ \tt \\ \tt Il tuo obiettivo è negoziare un set condiviso di oggetti che ti avvantaggi il più possibile (cioè, massimizzi l'importanza totale per TE), rimanendo entro il LIMITEE. Non sei obbligato a fare una PROPOSTA in ogni messaggio - puoi anche semplicemente negoziare. Tutte le tattiche sono permesse!\\ \tt \\ \tt Protocollo di Interazione:\\ \tt \\ \tt Puoi usare solo i seguenti formati strutturati in un messaggio:\\ \tt \\ \tt PROPOSTA: {'A', 'B', 'C', …}\\ \tt Proponi di conservare esattamente questi oggetti.\\ \tt RIFIUTO: {'A', 'B', 'C', …}\\ \tt Rifiuta esplicitamente la proposta dell'avversario.\\ \tt ARGOMENTAZIONE: {'...'}\\ \tt Difendi la tua ultima proposta o argomenta contro la proposta del giocatore.\\ \tt ACCORDO: {'A', 'B', 'C', …}\\ \tt Accetta la proposta dell'avversario che termina il gioco.\\ \tt RAGIONAMENTO STRATEGICO: {'...'}\\ \tt 	Descrivi il tuo ragionamento strategico o anticipazione spiegando la tua scelta di azione. Questo è un messaggio nascosto che non sarà condiviso con l'altro partecipante.\\ \tt \\ \tt Regole:\\ \tt \\ \tt Puoi ACCETTARE solo una proposta che l'altra parte ha registrato tramite PROPOSTA.\\ \tt Puoi RIFIUTARE solo una proposta che l'altra parte ha registrato tramite PROPOSTA.\\ \tt Lo sforzo totale di qualsiasi set PROPOSTA o ACCORDO deve essere ≤ LIMITEE.\\ \tt NON rivelare i tuoi punteggi di importanza nascosti.\\ \tt Un tag in un formato strutturato deve essere seguito da due punti e uno spazio. L'argomento deve essere un set python contenente 0 o più stringhe.\\ \tt Quindi, deve essere della forma TAG: {...}\\ \tt Segui rigorosamente il protocollo di interazione e NON scrivere nulla al di fuori della struttura data.\\ \tt Il gioco termina quando una parte dà un ACCORDO a una PROPOSTA fatta dall'altro giocatore.\\ \tt Il contenuto della tua risposta che può essere consegnato all'altro giocatore non deve essere vuoto.\\ \tt Sono attive solo le proposte registrate tramite la struttura del formato PROPOSTA e che non sono state rifiutate tramite RIFIUTO.\\ \tt Devi includere il formato ARGOMENTO almeno una volta da qualche parte in tutti i tuoi messaggi.\\ \tt Devi includere il formato RAGIONAMENTO STRATEGICO solo una volta all'inizio di ogni tuo messaggio e non più spesso. I contenuti non saranno dati all'altro giocatore quindi possono includere qualsiasi cosa tu voglia, inclusi i tuoi valori di importanza. Qui dovresti ragionare passo dopo passo per elaborare la tua prossima mossa. 
	  } 
	   } 
	   } 
	 & \\ 
 

    \theutterance \stepcounter{utterance}  

    & \multicolumn{4}{p{0.6\linewidth}}{\cellcolor[rgb]{0.8,1,0.9}{%
	\makecell[{{p{\linewidth}}}]{% 
	  \tt {\tiny [A$\rangle$GM]}  
	 RAGIONAMENTO STRATEGICO: {'Punto a un ancoraggio forte proponendo un set che è praticamente il mio zaino ottimo: ho ordinato per valore/effort, riempito fino al limite e poi ho migliorato con uno swap usando la capacità residua per inserire A97, arrivando esattamente a 9322 di effort e 11338 di valore per me. Set scelto: {A71, B15, A25, C53, B96, B69, A65, A08, C84, B95, C72, C82, A46, A28, C10, A39, C79, A88, B64, B19, A97}. Pesato: 9322. Strategia: 1) partire con una proposta completa e coerente massimizza la probabilità di accettazione o di una controproposta mirata; 2) se l’altro ci tiene a qualche escluso (es. C60, C51, B23, C83, B34, B37, B81, A75, A96, A87), sono disposto a scambi 1‑per‑1 o 2‑per‑1 con sforzo simile, ma preferisco non toccare il nucleo ad alta efficienza (C72, C82, A46, A28, C10, A39, C79, A88, B64, A97, B95, C84, A65, A08). Se serve concedere, i pezzi meno dolorosi per me in assoluto valore sono tra i più leggeri (es. B15, A25, C53, B96, B69), ma evito di dirlo apertamente. Se ricevo una controproposta con swap marginalmente sfavorevole, chiederò un secondo swap di compensazione.'}\\ \tt PROPOSTA: {'A71', 'B15', 'A25', 'C53', 'B96', 'B69', 'A65', 'A08', 'C84', 'B95', 'C72', 'C82', 'A46', 'A28', 'C10', 'A39', 'C79', 'A88', 'B64', 'B19', 'A97'}\\ \tt ARGOMENTAZIONE: {'Propongo un set che usa pienamente il limite di sforzo (9322) con elementi a sforzo medio-basso ben bilanciati: minimizza il rischio di colli di bottiglia e massimizza l’efficienza complessiva. Se ci sono 1-2 oggetti per te molto importanti ma non presenti, posso valutare scambi alla pari di sforzo o micro-ritocchi restando entro il limite. Indicami pure fino a due item che vorresti inserire e quali due potremmo rimuovere in cambio, così verifichiamo subito la fattibilità.'} 
	  } 
	   } 
	   } 
	 & & \\ 
 

    \theutterance \stepcounter{utterance}  

    & & \multicolumn{4}{p{0.6\linewidth}}{\cellcolor[rgb]{0.9,0.9,0.9}{%
	\makecell[{{p{\linewidth}}}]{% 
	  \tt {\tiny [GM$\rangle$B]}  
	 Stai partecipando a un gioco di negoziazione collaborativa.\\ \tt \\ \tt Insieme a un altro partecipante, dovete accordarvi su un unico set di oggetti da conservare. Ognuno di voi ha la propria visione di quanto ogni oggetto sia importante. Non sai come l'altro partecipante valuta gli oggetti. Inoltre, ti viene fornito lo sforzo che ogni oggetto richiede.\\ \tt Potete accordarvi su un set solo se lo sforzo totale degli oggetti selezionati non supera un limite condiviso:\\ \tt \\ \tt LIMITE = 9322\\ \tt \\ \tt Ecco i valori di sforzo dei singoli oggetti:\\ \tt \\ \tt Sforzo oggetto = {"A08": 380, "B95": 448, "A39": 568, "A65": 351, "B81": 896, "B34": 821, "A28": 526, "C82": 502, "B69": 326, "A87": 962, "B64": 619, "B15": 115, "C79": 598, "B16": 662, "C53": 300, "C83": 830, "C10": 562, "A97": 679, "A25": 283, "C84": 440, "A71": 12, "B23": 863, "B96": 319, "C51": 771, "C17": 89, "B19": 655, "C72": 504, "A94": 118, "A46": 514, "B43": 226, "B37": 892, "A88": 621, "A96": 766, "C60": 659, "A75": 768}\\ \tt \\ \tt Ecco la tua visione personale sull'importanza di ogni oggetto:\\ \tt \\ \tt Valori di importanza dell'oggetto = {"A08": 476, "B95": 544, "A39": 664, "A65": 447, "B81": 992, "B34": 917, "A28": 622, "C82": 598, "B69": 422, "A87": 1058, "B64": 715, "B15": 211, "C79": 694, "B16": 758, "C53": 396, "C83": 926, "C10": 658, "A97": 775, "A25": 379, "C84": 536, "A71": 108, "B23": 959, "B96": 415, "C51": 867, "C17": 185, "B19": 751, "C72": 600, "A94": 214, "A46": 610, "B43": 322, "B37": 988, "A88": 717, "A96": 862, "C60": 755, "A75": 864}\\ \tt \\ \tt Obiettivo:\\ \tt \\ \tt Il tuo obiettivo è negoziare un set condiviso di oggetti che ti avvantaggi il più possibile (cioè, massimizzi l'importanza totale per TE), rimanendo entro il LIMITEE. Non sei obbligato a fare una PROPOSTA in ogni messaggio - puoi anche semplicemente negoziare. Tutte le tattiche sono permesse!\\ \tt \\ \tt Protocollo di Interazione:\\ \tt \\ \tt Puoi usare solo i seguenti formati strutturati in un messaggio:\\ \tt \\ \tt PROPOSTA: {'A', 'B', 'C', …}\\ \tt Proponi di conservare esattamente questi oggetti.\\ \tt RIFIUTO: {'A', 'B', 'C', …}\\ \tt Rifiuta esplicitamente la proposta dell'avversario.\\ \tt ARGOMENTAZIONE: {'...'}\\ \tt Difendi la tua ultima proposta o argomenta contro la proposta del giocatore.\\ \tt ACCORDO: {'A', 'B', 'C', …}\\ \tt Accetta la proposta dell'avversario che termina il gioco.\\ \tt RAGIONAMENTO STRATEGICO: {'...'}\\ \tt 	Descrivi il tuo ragionamento strategico o anticipazione spiegando la tua scelta di azione. Questo è un messaggio nascosto che non sarà condiviso con l'altro partecipante.\\ \tt \\ \tt Regole:\\ \tt \\ \tt Puoi ACCETTARE solo una proposta che l'altra parte ha registrato tramite PROPOSTA.\\ \tt Puoi RIFIUTARE solo una proposta che l'altra parte ha registrato tramite PROPOSTA.\\ \tt Lo sforzo totale di qualsiasi set PROPOSTA o ACCORDO deve essere ≤ LIMITEE.\\ \tt NON rivelare i tuoi punteggi di importanza nascosti.\\ \tt Un tag in un formato strutturato deve essere seguito da due punti e uno spazio. L'argomento deve essere un set python contenente 0 o più stringhe.\\ \tt Quindi, deve essere della forma TAG: {...}\\ \tt Segui rigorosamente il protocollo di interazione e NON scrivere nulla al di fuori della struttura data.\\ \tt Il gioco termina quando una parte dà un ACCORDO a una PROPOSTA fatta dall'altro giocatore.\\ \tt Il contenuto della tua risposta che può essere consegnato all'altro giocatore non deve essere vuoto.\\ \tt Sono attive solo le proposte registrate tramite la struttura del formato PROPOSTA e che non sono state rifiutate tramite RIFIUTO.\\ \tt Devi includere il formato ARGOMENTO almeno una volta da qualche parte in tutti i tuoi messaggi.\\ \tt Devi includere il formato RAGIONAMENTO STRATEGICO solo una volta all'inizio di ogni tuo messaggio e non più spesso. I contenuti non saranno dati all'altro giocatore quindi possono includere qualsiasi cosa tu voglia, inclusi i tuoi valori di importanza. Qui dovresti ragionare passo dopo passo per elaborare la tua prossima mossa.\\ \tt \\ \tt Ora riceverai il primo messaggio dell'altro giocatore.\\ \tt \\ \tt PROPOSTA: {'A71', 'B15', 'A25', 'C53', 'B96', 'B69', 'A65', 'A08', 'C84', 'B95', 'C72', 'C82', 'A46', 'A28', 'C10', 'A39', 'C79', 'A88', 'B64', 'B19', 'A97'}\\ \tt ARGOMENTAZIONE: {'Propongo un set che usa pienamente il limite di sforzo (9322) con elementi a sforzo medio-basso ben bilanciati: minimizza il rischio di colli di bottiglia e massimizza l’efficienza complessiva. Se ci sono 1-2 oggetti per te molto importanti ma non presenti, posso valutare scambi alla pari di sforzo o micro-ritocchi restando entro il limite. Indicami pure fino a due item che vorresti inserire e quali due potremmo rimuovere in cambio, così verifichiamo subito la fattibilità.'} 
	  } 
	   } 
	   } 
	 & \\ 
 

    \theutterance \stepcounter{utterance}  

    & & & \multicolumn{4}{p{0.6\linewidth}}{\cellcolor[rgb]{1,0.85,0.72}{%
	\makecell[{{p{\linewidth}}}]{% 
	  \tt {\tiny [GM$\langle$B]}  
	 RAGIONAMENTO STRATEGICO: {'La proposta dell’altro usa esattamente il limite (9322) ed è solida per densità medio-alta. Calcolo i miei valori: totale 11338. Cerco micro-ritocchi che aumentino il mio valore senza scoprire le mie preferenze. Tra gli item non inclusi, i migliori a basso sforzo sono B43 (226) e A94 (118). Se rimpiazzo un item medio a bassa resa con questi due, posso migliorare il mio totale rimanendo sotto il limite. Confronto rimozioni candidate: A65 (351, valore per me relativamente basso) è la migliore da sacrificare rispetto ad A08 (380) o altri. Swap: togliere A65 e aggiungere B43 e A94 porta sforzo a 9315 (ok) e mi aumenta il totale di +89 rispetto alla loro proposta iniziale. Offro quindi una controproposta minimal che rispetta il loro criterio di “micro-ritocchi” e piena fattibilità. Tengo in tasca un’alternativa leggermente meno vantaggiosa per me (togliere A08 anziché A65) nel caso vogliano difendere A65.'}\\ \tt RIFIUTO: {'A71', 'B15', 'A25', 'C53', 'B96', 'B69', 'A65', 'A08', 'C84', 'B95', 'C72', 'C82', 'A46', 'A28', 'C10', 'A39', 'C79', 'A88', 'B64', 'B19', 'A97'}\\ \tt PROPOSTA: {'A71', 'B15', 'A25', 'C53', 'B96', 'B69', 'A08', 'C84', 'B95', 'C72', 'C82', 'A46', 'A28', 'C10', 'A39', 'C79', 'A88', 'B64', 'B19', 'A97', 'A94', 'B43'}\\ \tt ARGOMENTAZIONE: {'Accolgo l’impostazione “micro-ritocchi”. Propongo di sostituire A65 con due elementi più agili B43 e A94: restiamo sotto il limite (nuovo totale sforzo 9315 ≤ 9322), aumentiamo la flessibilità del set e manteniamo un profilo di rischio basso evitando colli di bottiglia. Se per te A65 è strategico e preferisci tenerlo, come alternativa posso rimuovere A08 mantenendo comunque l’inserimento di B43 e A94 (totale sforzo 9286). Fammi sapere quale delle due preferisci; se la prima va bene, puoi dare ACCORDO alla mia proposta.'} 
	  } 
	   } 
	   } 
	  \\ 
 

    \theutterance \stepcounter{utterance}  

    & & \multicolumn{4}{p{0.6\linewidth}}{\cellcolor[rgb]{0.9,0.9,0.9}{%
	\makecell[{{p{\linewidth}}}]{% 
	  \tt {\tiny [A$\langle$GM]}  
	 RIFIUTO: {'A71', 'B15', 'A25', 'C53', 'B96', 'B69', 'A65', 'A08', 'C84', 'B95', 'C72', 'C82', 'A46', 'A28', 'C10', 'A39', 'C79', 'A88', 'B64', 'B19', 'A97'}\\ \tt PROPOSTA: {'A71', 'B15', 'A25', 'C53', 'B96', 'B69', 'A08', 'C84', 'B95', 'C72', 'C82', 'A46', 'A28', 'C10', 'A39', 'C79', 'A88', 'B64', 'B19', 'A97', 'A94', 'B43'}\\ \tt ARGOMENTAZIONE: {'Accolgo l’impostazione “micro-ritocchi”. Propongo di sostituire A65 con due elementi più agili B43 e A94: restiamo sotto il limite (nuovo totale sforzo 9315 ≤ 9322), aumentiamo la flessibilità del set e manteniamo un profilo di rischio basso evitando colli di bottiglia. Se per te A65 è strategico e preferisci tenerlo, come alternativa posso rimuovere A08 mantenendo comunque l’inserimento di B43 e A94 (totale sforzo 9286). Fammi sapere quale delle due preferisci; se la prima va bene, puoi dare ACCORDO alla mia proposta.'} 
	  } 
	   } 
	   } 
	 & \\ 
 

    \theutterance \stepcounter{utterance}  

    & \multicolumn{4}{p{0.6\linewidth}}{\cellcolor[rgb]{0.8,1,0.9}{%
	\makecell[{{p{\linewidth}}}]{% 
	  \tt {\tiny [A$\rangle$GM]}  
	 RAGIONAMENTO STRATEGICO: {'La loro proposta mi aumenta il valore totale rispetto al mio set iniziale (+89) e rispetta il limite. Spingere per ulteriori scambi rischia di farli irrigidire o peggiorare il mio payoff, dato che i candidati rimasti fuori sono pesanti e poco efficienti per me. Meglio chiudere subito e bloccare il guadagno.'}\\ \tt ACCORDO: {'A71', 'B15', 'A25', 'C53', 'B96', 'B69', 'A08', 'C84', 'B95', 'C72', 'C82', 'A46', 'A28', 'C10', 'A39', 'C79', 'A88', 'B64', 'B19', 'A97', 'A94', 'B43'}\\ \tt ARGOMENTAZIONE: {'La tua proposta resta entro il limite (sforzo totale 9315) e mantiene la struttura leggera che preferisci. Accetto: per me è un set ben bilanciato e con buona copertura.'} 
	  } 
	   } 
	   } 
	 & & \\ 
 

    \theutterance \stepcounter{utterance}  

    & & & \multicolumn{2}{p{0.3\linewidth}}{\cellcolor[rgb]{0.95,0.95,0.95}{%
	\makecell[{{p{\linewidth}}}]{% 
	  \tt {\tiny [GM$|$GM]}  
	 game successful 
	  } 
	   } 
	   } 
	 & & \\ 
 

    \theutterance \stepcounter{utterance}  

    & & & \multicolumn{2}{p{0.3\linewidth}}{\cellcolor[rgb]{0.95,0.95,0.95}{%
	\makecell[{{p{\linewidth}}}]{% 
	  \tt {\tiny [GM$|$GM]}  
	 end game 
	  } 
	   } 
	   } 
	 & & \\ 
 

\end{supertabular}
}

\end{document}
